\section{Limitations}
\label{sec:limitations}

\paragraph{Typed Transformation Assumption.}

Our formulation assumes that tools can be modeled as typed transformations between entity types.
This assumption holds naturally for infrastructure APIs, platform operations, and many data-processing workflows, where tools consume one resource and produce another.
It is less natural for imperative or side-effect operations (e.g., \texttt{send\_email} or \texttt{restart\_service}) that do not naturally expose typed outputs.
Such operations can be accommodated using terminal entity types (e.g., \texttt{Action}), but doing so requires additional modeling decisions during graph construction.

\paragraph{Sequential Workflows.}

The current formulation recovers a single shortest path and therefore models sequential workflows only.
Branching execution, parallel composition, and conditional control flow would require extending the graph representation beyond simple paths.

\paragraph{Uniform Edge Costs.}

Breadth-first search treats all edges equally.
In practice, a write operation may appear on a shorter path than a semantically preferable read-only sequence.
Weighted graph search using operation type or execution cost is a natural extension.

\paragraph{Entity Type System Quality.}

The effectiveness of graph routing depends on the granularity of the entity type system.
Our experiments use manually constructed type systems; automatically deriving typed composition graphs from API specifications (e.g., OpenAPI schemas) remains an open direction.

\paragraph{Single-Intent Queries.}

Our evaluation focuses on single-objective queries.
Compound requests naturally correspond to multiple independent graph paths and would require an intent decomposition stage before graph search.
